\documentclass[12pt]{article}
\usepackage[spanish]{babel}
\usepackage[utf8]{inputenc}
\usepackage[T1]{fontenc}
\usepackage{amsmath, amssymb}
\usepackage{geometry}
\usepackage{enumitem}
\usepackage{needspace}
\usepackage{tikz}
\usetikzlibrary{babel}

\geometry{letterpaper, margin=2cm}
\setlength{\parindent}{0pt}
\setlength{\parskip}{6pt}

\newcommand{\puntografico}{0.35}
\newcommand{\opcionalgebra}[2]{\item[#1)] $#2$}
\newcommand{\grafcirc}[5]{%
\begin{tikzpicture}[scale=0.078]
  \path[use as bounding box] (-32,-15) rectangle (43,34);
  \draw[<->, thick] (-30,0) -- (36,0);
  \draw[<->, thick] (0,-12) -- (0,31);
  \node[anchor=west] at (37,-1.2) {$x$ (este)};
  \node[anchor=south east] at (7,32.5) {$y$ (norte)};
  \fill (0,0) circle (\puntografico) node[below left] {$P$};
  \draw[dashed] (#1,0) -- (#1,#2);
  \draw[dashed] (0,#2) -- (#1,#2);
  \draw[thick] (#1,#2) circle (#3);
  \fill (#1,#2) circle (\puntografico);
  \node[anchor=south, inner sep=1pt, yshift=1pt] at (#1,#2) {$#4$};
  \node[anchor=south west, inner sep=0.5pt] at (#1+#3*0.72,#2+#3*0.78) {$C$};
  \node[anchor=north] at (#1,0) {$#5$};
  \node[left] at (0,#2) {$#2$};
\end{tikzpicture}}

\begin{document}

\begin{center}
  {\Large \textbf{Borrador de items}}\\
  {\large Bloque 1. Geometria - Afirmacion 1}\\[4pt]
  Prueba Nacional Estandarizada de Matematica, Secundaria 2026
\end{center}

\section*{Afirmacion}
Resuelve problemas relacionados con la representacion de circunferencias de manera analitica o grafica.

\section*{Items propuestos}

\begin{enumerate}[label=\textbf{\arabic*.}, itemsep=1.05cm]

\Needspace{0.92\textheight}
\item
El alcance maximo de la senal inalambrica que emite un dispositivo es de $8$ m a su alrededor. La ubicacion de ese dispositivo ($M$) es $20$ m al este y $12$ m al norte de la ubicacion de una casa ($P$), la cual se considera como origen.

\textquestiondown Cual de las siguientes representaciones graficas, en la que las unidades estan en metros, corresponde a la circunferencia del alcance maximo de la senal que emite ese dispositivo?

\vspace{2\baselineskip}

\begin{center}
\setlength{\tabcolsep}{0pt}
\renewcommand{\arraystretch}{1}
\begin{tabular}{@{}r@{\hspace{0.45cm}}c@{\hspace{1.7cm}}r@{\hspace{0.45cm}}c@{}}
\raisebox{1.65cm}{\textbf{A)}} & \grafcirc{20}{12}{8}{M}{20} &
\raisebox{1.65cm}{\textbf{B)}} & \grafcirc{12}{20}{8}{M}{12} \\[0.85cm]
\raisebox{1.65cm}{\textbf{C)}} & \grafcirc{20}{12}{4}{M}{20} &
\raisebox{1.65cm}{\textbf{D)}} & \grafcirc{-20}{12}{8}{M}{\llap{-}20}
\end{tabular}
\end{center}

\Needspace{0.76\textheight}
\item
La siguiente representacion grafica, en la que las unidades estan en centimetros, muestra la ubicacion del centro ($P$) de la base circular de una pecera y la circunferencia correspondiente al borde de esa base.

\begin{center}
\begin{tikzpicture}[scale=0.18]
  \path[use as bounding box] (-8,-8) rectangle (28,18);
  \draw[<->, thick] (-5,0) -- (25,0);
  \draw[<->, thick] (0,-5) -- (0,16);
  \node[anchor=west] at (26,-0.8) {$x$};
  \node[anchor=south east] at (3.5,16.6) {$y$};
  \draw[dashed] (10,0) -- (10,6);
  \draw[dashed] (0,6) -- (10,6);
  \draw[thick] (10,6) circle (6);
  \fill (10,6) circle (\puntografico) node[above left] {$P$};
  \node[anchor=north] at (10,0) {$10$};
  \node[left] at (0,6) {$6$};
\end{tikzpicture}
\end{center}

De acuerdo con la informacion anterior, \textquestiondown cual es la representacion algebraica, en la que las unidades estan en centimetros, de la circunferencia correspondiente al borde de la base de la pecera?

\begin{enumerate}[label=\Alph*)]
  \opcionalgebra{A}{(x+10)^2+(y+6)^2=36}
  \opcionalgebra{B}{(x-10)^2+(y-6)^2=36}
  \opcionalgebra{C}{(x-6)^2+(y-10)^2=36}
  \opcionalgebra{D}{(x-10)^2+(y-6)^2=6}
\end{enumerate}

\Needspace{0.88\textheight}
\item
Una antena de telecomunicaciones ($T$) emite una senal que permite a los telefonos celulares, ubicados a $13$ km o menos alrededor de esa antena, realizar y recibir llamadas. En una representacion grafica, la ubicacion de la antena corresponde al punto $T(5,12)$.

Durante una revision tecnica se registraron tres telefonos celulares ubicados en los puntos $R(10,24)$, $S(18,12)$ y $U(17,18)$.

La siguiente representacion grafica, en la que las unidades estan en kilometros, muestra la ubicacion de la antena y la circunferencia correspondiente al alcance maximo de la senal.

\begin{center}
\begin{tikzpicture}[scale=0.16]
  \path[use as bounding box] (-6,-5) rectangle (27,31);
  \draw[<->, thick] (-4,0) -- (25,0);
  \draw[<->, thick] (0,-3) -- (0,29);
  \node[anchor=west] at (26,-0.8) {$x$ (este)};
  \node[anchor=south east] at (5.5,30) {$y$ (norte)};
  \draw[thick] (5,12) circle (13);
  \draw[dashed] (5,0) -- (5,12);
  \draw[dashed] (0,12) -- (5,12);
  \fill (5,12) circle (\puntografico);
  \node[anchor=south east, inner sep=2pt, xshift=-2pt, yshift=2pt] at (5,12) {$T$};
  \node[anchor=north] at (5,0) {$5$};
  \node[left] at (0,12) {$12$};
\end{tikzpicture}
\end{center}

De acuerdo con la informacion anterior, \textquestiondown en cual de esos telefonos se pueden realizar y recibir llamadas por medio de la senal de esa antena?

\begin{enumerate}[label=\Alph*)]
  \item Solo en $R$.
  \item Solo en $S$.
  \item En $R$ y en $S$.
  \item En $R$, en $S$ y en $U$.
\end{enumerate}

\Needspace{0.88\textheight}
\item
En una finca se desea construir un estanque circular para peces. El centro del estanque ($C$) se ubicara a $9$ m al oeste y $4$ m al sur de un arbol ($A$), el cual se considera como origen. Ademas, el borde del estanque debe quedar a $7$ m del centro.

La siguiente representacion grafica, en la que las unidades estan en metros, muestra la ubicacion del arbol, del centro del estanque y de la circunferencia correspondiente al borde del estanque.

\begin{center}
\begin{tikzpicture}[scale=0.18]
  \path[use as bounding box] (-22,-16) rectangle (13,10);
  \draw[<->, thick] (-20,0) -- (11,0);
  \draw[<->, thick] (0,-14) -- (0,8);
  \node[anchor=west] at (12,-0.8) {$x$ (este)};
  \node[anchor=south east] at (4,9) {$y$ (norte)};
  \fill (0,0) circle (\puntografico) node[below right] {$A$};
  \draw[dashed] (-9,0) -- (-9,-4);
  \draw[dashed] (0,-4) -- (-9,-4);
  \draw[thick] (-9,-4) circle (7);
  \fill (-9,-4) circle (\puntografico) node[anchor=east, inner sep=2pt] {$C$};
  \node[anchor=north] at (-9,0) {$\llap{-}9$};
  \node[right] at (0,-4) {$-4$};
\end{tikzpicture}
\end{center}

\textquestiondown Cual de las siguientes representaciones algebraicas, en la que las unidades estan en metros, corresponde a la circunferencia que representa el borde de ese estanque?

\begin{enumerate}[label=\Alph*)]
  \opcionalgebra{A}{(x-9)^2+(y-4)^2=49}
  \opcionalgebra{B}{(x+9)^2+(y+4)^2=49}
  \opcionalgebra{C}{(x+9)^2+(y+4)^2=7}
  \opcionalgebra{D}{(x-4)^2+(y-9)^2=49}
\end{enumerate}

\end{enumerate}

\end{document}
